\documentclass[12pt,a4paper]{article}

\setlength{\topmargin}{0.0in}
\setlength{\oddsidemargin}{0.33in}
\setlength{\textheight}{9.0in}
\setlength{\textwidth}{6.0in}
\renewcommand{\baselinestretch}{1.25}



\usepackage{hyperref}
% \usepackage[utf8]{inputenc}
% \usepackage[T1]{fontenc}
\usepackage{float}
\usepackage{enumitem}
\usepackage{amsmath}
\usepackage{amssymb}
\usepackage{tikz}
\usetikzlibrary{decorations.pathreplacing}
\usepackage{amsfonts}
\usepackage{graphicx}
\usepackage{subcaption}
\usepackage{caption}
\newcommand{\dv}[2]{\frac{\mathrm{d} #1}{\mathrm{d} #2}}
\usepackage{caption}
\captionsetup[figure]{font=small} 
% \captionsetup[table]{font=small}
\usepackage[capitalize,nameinlink]{cleveref}
\usepackage{mathtools}
\newcommand{\defeq}{\vcentcolon=}
\newcommand{\eqdef}{=\vcentcolon}

% Define environments
\newtheorem{theorem}{Theorem}
\newtheorem{definition}{Definition}
\newtheorem{lemma}{Lemma}
\newtheorem{corollary}{Corollary}



\title{Meeting Notes August 11}
\author{}
\date{}
% \author{Maggie McCarthy}
% \date{May 2026}

\begin{document}

\maketitle


\section{Code Updates}

\begin{itemize}
\item Everything is running in parallel (4 ranks) and saving data each sweep. 
\item I believe that the Dirac Poisson run looks good. I have all the data saved so that I can alter the plots as needed.
\item Maxwell is challenging. I have access to the Prague computer and a better departmental computer. Even then, a standard run is a challenge. Once we hit the second sweep, we are hitting max DOF (20,000) at almost every test point. 
\item However, there is some potential good news. I think there is a way to massively speed things up. I have worked this out below. I have also done some little numerical tests to check this. I am running this code now and I expect it to be finished today. If so, the whole run, with max dof set to 100,000, would take about 24 hours, which is a dramatic speed up.
\item However, it is still a challenge to get results that look as good as the Dirac Poisson. What is the conclusion here? 
\end{itemize}



\section{Writing Updates}

\subsection{NOT Done: (Draft)}
\begin{itemize}
    \item Folding with FEM.
    \item Results section.
    \item Conclusions, Abstract, Introduction.
\end{itemize}

\subsection{Questions}
\begin{itemize}
\item What time on Aug 24? What is needed by Aug 24? How complete? How edited? 
\item Worried about being over page count by Aug 24.
\item  Is it ok to present Poisson and Maxwell? Do I need other examples? I am challenged with respect to time and space in the write-up. Umberto's examples were
    \begin{itemize}
        \item 1D Poisson problem,
        \item 1D advection diffusion,
        \item 2D Maxwell on a square,
        \item 2D Maxwell on an L-shape,
        \item 1D Cossu and Chomaz operator,
        \item 1D Benilov, O'Brien and Sazonov operator,
        \item 2D Maxwell with Lie-Derivative. 
    \end{itemize}
\item What do we think the folding with FEM section will look like?
\item I am not quite following Umberto's domain issues - page 18.
\item Page 23 - do we need some sort of embedding for DWR?
\item How to tackle the discrete bottom of spectrum discussion (page 29). Do I need it in an Appendix?
\item Figure 3.2 - should I take out the orange curve?
\end{itemize}

\section{Meeting with Matt}

\begin{itemize}
    \item I will prepare some slides today. What should be in them exactly?
    \item What else needs to be prepped? 
\end{itemize}


\section{Speeding up Outer Loop (Self-Adjoint Problems)}

\subsection{Matt's Folded Operator}

Let $z = x + iy.$ Matt's folded operator is $F(z) \defeq (A-zI)^*(A-zI).$ Assume that $A^*=A.$ Then,
$F(z) = (A-\overline{z}I)(A-zI).$

Next, plug in $z$ and expand
\begin{align}
    F(z) &= (A-xI + iyI)(A-xI-iyI)\\
    &=(A-zI)^2 + y^2 I.
\end{align}

Furthermore, for the injection modulus, we have
\begin{align}
    \|(A-zI) u \|^2 &= \| (A-zI) u - iyu \|^2 \\
    &= \| (A-xI)u \|^2 + y^2 \|u\|^2 + 2 \Re\{ (A-xI)u, -iyu ) \} \\
    &= \| (A-xI)u \|^2 + y^2 \|u\|^2.
\end{align}
Taking the infimum over unit vectors, we obtain
\begin{equation}
    \gamma(z, A)^2 = \gamma(x, A)^2 + y^2. 
\end{equation}

\subsection{FEM Eigenproblem}

Next, we consider $B_z \defeq (A-zM)^*M^{-1}(A-zM) = (A-\overline{z}M)M^{-1}(A-zM).$ Expanding, we have 
\begin{align}
    B_z &= AM^{-1}A - (z + \overline z) A + |z|^2 M \\
    &= AM^{-1} A - 2xA + (x^2 + y^2) M.
\end{align}
Notice that 
\begin{equation}
    B_x = (A-xM)M^{-1}(A-xM) = AM^{-1}A - 2xA + x^2 M.
\end{equation}
Thus, 
\begin{equation}
    B_z = B_x + y^2 M.
\end{equation}
Now, define the sesquilinear forms 
\begin{align}
    a_z(u_h, v_h) &= \langle B_z u_h, v_h \rangle \\
    &= \langle B_x u_h, v_h \rangle + y^2 \langle M u_h, v_h \rangle  \\
    &= a_x(u_h, v_h) + y^2 m(u_h, v_h).
\end{align}
Thus, the eigenproblem 
\begin{equation}
    a_z(u_h, v_h) = \lambda_h(z) m(u_h, v_h),
\end{equation}
can be written as 
\begin{equation}
    a_x(u_h, v_h) = (\lambda_h(z) - y^2) m(u_h, v_h) = \lambda_h(x) m(u_h, v_h).
\end{equation}
In other words, for each distinct $x,$ we can solve $a_x(u_h, v_h) = \lambda_h(x) m(u, v)$ and find $\lambda_h(z)$ by adding $y^2.$ Furthermore, the eigenfunctions carry over as well. That is, $E_h(z) = E_h(x),$ where $E_h(\cdot)$ denotes the eigenspace for $\lambda_h(\cdot).$

With this, we can conclude that for $z$ values that share the same real part, Matt's bound on the distance to the spectrum depends only on a shift by the square of the imaginary part:
\begin{equation}
    \Phi_n(z, A) = \sqrt{\lambda_h(z)} = \sqrt{\lambda_h(x) + y^2}.
\end{equation}


\subsection{DWR Error Estimate}

Notice that the imaginary part of $z$ does not affect the primal residual:
\begin{align}
    \rho_z(v) &\defeq a_z(u_h, v) - \lambda_h(z) m(u_h, v) \\
    &= a_x(u_h, v) + y^2 m(u_h, v) - \lambda_h(z) m(u_h, v) \\
    &= a_x(u_h, v) - (\lambda_h(z) - y^2) m(u_h, v) \\
    &= a_x(u_h, v) - \lambda_h(x) m(u_h, v) \\
    &= \rho_x(v).  
\end{align}
Since the residual only depends on $x,$ and the eigenfunctions only depend on $x,$ the DWR error estimate depends only on $x$:
\begin{align}
    \lambda(z) - \lambda_h(z) &\approx \frac{\rho_z(u_h^+(z) - u_h(z))}{1 - 1/2 \| u_h^+(z) - u_h(z) \|_m^2 }\\
    &= \frac{\rho_z(u_h^+(x) - u_h(x))}{1 - 1/2 \| u_h^+(x) - u_h(x) \|_m^2 }.
\end{align}


\subsection{What does this mean?}

If correct, this means we can drastically speed up our computations. For each $x,$ we do the expensive adaptive eigensolves. Then, for each `column' (every $z$ that shares $x$) in the complex plane, we compute the needed $\Phi_n(z, A)$ using the shift by $y^2.$ The error estimate computed at $x$ carries over to every other point in that column.



% \section{The Functional Analysis Issues}


% \subsection{Matt's Formulations:}

% Matt considers a closed and densely defined operator $T : D(T) \subset X \rightarrow X,$ where $X$ is a Hilbert space. In doing so, he defines the spectrum and pseudospectrum using $T-zI.$ He folds using $(T-zI)^*(T-zI)$ where $(T-zI)^*$ is the Hilbert adjoint. Furthermore, he defines the injection modulus as follows:
% \begin{equation}
%     \sigma_{\inf}(T-zI) = \inf_{u \in D(T), u \ne 0} \frac{ \| (T-zI)u \|_X} {\|u\|_X}.
% \end{equation}

% \subsection{The Issue}

% When we work on the Maxwell problem and/or the Poisson problem, we consider an operator $A : V \rightarrow V^*.$ Thus, it does not make sense to write $A-zI.$ What is $I$?

% Furthermore, even if we replace $I$ with the correct identity map (say a Riesz map), it does not make sense to fold $(A-zM)^*(A-zM)$ because there is a space mismatch:
% \begin{equation}
%     A-ZM : V \rightarrow V^*
% \end{equation}
% \begin{equation}
%     (A-ZM)^* : V \rightarrow V^*,
% \end{equation}
% where $(A-zM)^*$ is the Banach adjoint. 

% As a result, we need to rethink the functional analytic setting for Matt's algorithm so that it works for the problems we want to work with.

% \subsection{Why Our Previous Solutions Failed}

% We considered $A : V \rightarrow V^*$ and a special $M: V \rightarrow V^*$ defined via $\langle Mu, v\rangle = (u, v)_{L^2}.$ This is not surjective! So we could not honestly invert $M^{-1}:V^* \rightarrow V.$ Also, we could not connect $A$ back to $T$ without a Riesz map. 

% \subsection{Current Solution}

% Let $V = H_0^1 \times H^1$ for Poisson (dirac folding) or $V = H_{\text{curl}} \times H_{\text{curl}} $ for Maxwell (2D). Set $X = L^2 \times L^2.$ We know $V \hookrightarrow X.$

% The key is that we use first-order formulations of the two operators. For example, for Poisson, we use $TU = (-i \sigma', -i u')^T$ where $U = (u, \sigma)^T.$ Thus, our operator $A$ maps from V to $X^*.$ To see this, again consider Poisson: since $u \in H_0^1$ and $\sigma \in H^1,$ $TU \in L^2.$ Thus, we can define $\langle AU, V \rangle = (TU, V)_{L^2}$ where $A : V \rightarrow X^*.$

% With this, we can consider $M$ as the inverse Riesz map:
% \begin{equation}
%     M: X \rightarrow X^*,\qquad \langle MU, V\rangle_{X^*, X} = (U, V)_X.
% \end{equation}
% The Riesz representation theorem gives the related Riesz map $M^{-1}: X^* \rightarrow X.$ 

% With this, we can write $A-zM: V \rightarrow X^*$ because $A : V \rightarrow X^*,$ $M: X \rightarrow X^*,$ and any $u \in V$ is also in $X.$

% Furthermore, we can compose $(A-zM)^* M^{-1} (A-zM)$ because $(A-zM): V \rightarrow X^*,$ $M^{-1}: X^* \rightarrow X,$ and $(A-zM)^* : X \rightarrow V^*$ (banach adjoint). Thus, our new folded operator is 
% \begin{equation}
%     (A-zM)^* M^{-1} (A-zM) : V \rightarrow V^*.
% \end{equation}

% The KEY to this result is that we have a first-order operator to work with. Otherwise, we have $A: V \rightarrow V^*$ and not $A: V \rightarrow X^*$ and choosing the proper $M$ for folding was a disaster because we could not use the standard Riesz map from $X^*$ to $X.$ 


% \subsection{Connection to Matt's Formulation}

% Now, return to Matt's $T: V \rightarrow X$ which is assumed to be a first-order operator. We define 
% \begin{equation}
%     A : V \rightarrow X^*, \qquad \langle Au, v \rangle_{X^*, X} = (Tu, v)_{X}.
% \end{equation}
% Then, we can say $A = MT$ where $M:X \rightarrow X^*$ is the inverse Riesz map. Similarly, we can say $(A-zM) = M(T-zI).$

% In this way, we can define $\operatorname{Sp}(A; M) =\{z \in \mathbb{C} : A-zM \text{ is not boundedly invertible} \}$ which is equivalent to $\operatorname{Sp}(T)$ because $(T-zI)$ is boundedly invertible iff $A-zM$ is ($M$ is an isomorphism). 

% Furthermore, we can show the equivalence of the injection modulus. By the Riesz representation theorem, we have 
% $ \| Mf \|_{X^*} = \|f \|_{X}.$ Thus, 
% \begin{align}
%     \sigma_{\inf}(T-zI) &= \inf_{u \in V, u \ne 0} \frac{\|(T-zI) u \|_{X}}{\|u\|_X} \\\
%     &= \inf_{u \in V, u \ne 0} \frac{\|M (T-zI) u \|_{X^*}}{\|Mu\|_{X^*}} \\
%     &= \inf_{u \in V, u \ne 0} \frac{\|(A-zM) u \|_{X^*}}{\|u\|_X}.
% \end{align}
% Thus, $\sigma_{\inf}(T-zI) = \sigma_{\inf}((A-zM); M).$

% Furthermore,
% \begin{align}
%     \sigma_{\inf}^2(A-zM;M) &= \inf_{u \in V, u \ne 0} \frac{\|(A-zM) u \|_{X^*}^2}{\|u\|_X^2}\\
%     &= \inf_{u \in V, u \ne 0} \frac{ \langle (A-zM)u, M^{-1} (A-zM)u \rangle } {\langle Mu, u \rangle} \\
%     &= \inf_{u \in V, u \ne 0} \frac{ \langle (A-zM)^* M^{-1} (A-zM) u, u \rangle} {\langle Mu, u \rangle},
% \end{align}
% which is the folding we currently implement (or want to be implementing). 


% \subsection{Example: Poisson}

% We define $V = H_0^1 \times H^1$ and $X = L^2 \times L^2.$ We consider $T:V \rightarrow X$ defined by 
% \begin{equation}
%     T(u, \sigma) = (-i \sigma', -i u')^T. 
% \end{equation}

% We define $A:V \rightarrow X^*$ by 
% \begin{equation}
%     \langle A(u, \sigma), (v, \tau)\rangle = (T(u, \sigma), (v, \tau))_{X} = (-i \sigma', v)_{L^2} + (-iu', \tau)_{L^2}.
% \end{equation}
% Defining $M:X \rightarrow X^*$  by $\langle Mu, v \rangle = (u, v)$ we have $A = MT$ because, by definiton of the inverse Riesz map,
% \begin{equation}
%     \langle MTu, v\rangle_{X^*, X} = (Tu, v)_X = \langle Au, v \rangle\rangle_{X^*, X}. 
% \end{equation}

% Thus, we can use the $\sigma_{\inf}(A;M)$ to recover the desired pseudospectral information about $T.$ 


% \subsection{Still to Sort}
% \begin{itemize}
%     \item I need to check that the same result holds for $\sigma_{\inf}(T^* - \overline{z}I).$
%     \item Umberto and I need to prove the desired Gårding inequality so that we have a discrete spectrum. 
% \end{itemize}
% \subsection{Questions}
% \begin{itemize}
%     \item Dual or Antidual?
%     \item Is this the correct direction to go in? 
%     \item How should this be presented in the thesis? 
% \end{itemize}




\end{document}